\documentclass[a4paper,12pt]{article}

% =========================
% Pakete
% =========================
\usepackage[utf8]{inputenc}
\usepackage[ngerman]{babel}
\usepackage{lmodern}
\usepackage{geometry}
\geometry{margin=2.5cm}
\usepackage{setspace}
\onehalfspacing

% =========================
% Dokument
% =========================
\begin{document}

\title{Projektübersicht: Multimodaler Sprach-Assistent}
\author{}
\date{}
\maketitle

\section{Projektbeschreibung}

Das Ziel des Projekts ist die Entwicklung eines intelligenten, Python-basierten Chatbots, der als interaktive Schnittstelle zwischen Mensch und Maschine dient. Das System verarbeitet gesprochene Sprache und kombiniert dabei Methoden der Spracherkennung, natürlichen Sprachverarbeitung und visuellen Informationsdarstellung.

\section{Spracherkennung und Textverarbeitung}

Zunächst nimmt das System gesprochene Sprache über ein Mikrofon auf und wandelt diese mithilfe von Speech-to-Text-Technologien in geschriebenen Text um. Dieser transkribierte Text bildet die Grundlage für die weitere Verarbeitung.

Im nächsten Schritt erfolgt eine linguistische Analyse des Textes. Dabei wird der Satz in einzelne Bestandteile zerlegt (Tokenisierung). Anschließend werden grammatikalische Wortarten bestimmt (Part-of-Speech Tagging), wobei insbesondere Verben und Nomen im Fokus stehen.

\section{Informationsextraktion}

Nach der Analyse werden relevante Informationen aus dem Text extrahiert. Insbesondere werden Verben und konkrete Objekte (Nomen) identifiziert. Diese strukturierten Informationen dienen als Grundlage für die weitere Verarbeitung und Visualisierung.

\section{Visuelle Darstellung}

Die extrahierten Informationen werden genutzt, um eine visuelle Ausgabe zu erzeugen. Dabei gibt es zwei zentrale Funktionen:
\begin{itemize}
    \item \textbf{Textuelle Visualisierung:} Verben werden im ursprünglichen Text hervorgehoben (z.B. farblich markiert), um sprachliche Strukturen sichtbar zu machen.
    \item \textbf{Bildgenerierung:} Erkannte Objekte werden genutzt, um passende Bilder aus einer Bilddatenbank abzurufen oder mithilfe generativer KI zu erzeugen.
\end{itemize}

\section{Ziel des Projekts}

Das Projekt verfolgt das Ziel, verschiedene Kernbereiche der künstlichen Intelligenz in einem interaktiven System zu vereinen. Dazu gehören insbesondere Spracherkennung, Textanalyse und visuelle Darstellung.

Der entwickelte Chatbot dient als intuitive Benutzeroberfläche, die gesprochene Sprache in strukturierte und visuell aufbereitete Informationen transformiert. Dadurch wird die Abstraktion der Sprache für den Nutzer greifbar und verständlich gemacht.

\end{document}